% Render: python3 ~/.agents/skills/create-infographic/scripts/render_tikz.py docs/architecture-infographic.tex --workdir . --out-dir docs --engine xelatex --dpi 300 --svg --strict-warnings --strict-geometry --strict-composition --strict-svg --min-labelled-connectors 0.6 --expect-text "generative-art-demos"
\documentclass[tikz,border=0pt]{standalone}
\usepackage[UTF8,fontset=fandol]{ctex}
% Prefer Source Han Sans SC when installed (exact em-width glyph boxes);
% the ctex fandol set above remains the portable fallback.
\IfFontExistsTF{Source Han Sans SC}{%
  \setCJKmainfont{Source Han Sans SC}%
  \setCJKsansfont{Source Han Sans SC}%
}{}
\usepackage{tikz}
\usetikzlibrary{shapes.geometric, arrows.meta, positioning, calc, decorations.pathreplacing, decorations.pathmorphing, patterns, fit, backgrounds}
\usepackage{xcolor}

% Color palette matching the project
\definecolor{primary}{HTML}{ff6f91}
\definecolor{accent1}{HTML}{ffd166}
\definecolor{accent2}{HTML}{54d6c6}
\definecolor{accent3}{HTML}{7b9cff}
\definecolor{accent4}{HTML}{ba8cff}
\definecolor{darkbg}{HTML}{0f1e2d}
\definecolor{midbg}{HTML}{1a2f44}
\definecolor{lighttext}{HTML}{e8f4fc}
\definecolor{mutetext}{HTML}{89a6bb}
\definecolor{bordercol}{HTML}{2d4a63}
\definecolor{gridcol}{HTML}{1e3547}

\begin{document}
\begin{tikzpicture}[
  x=1pt, y=1pt,
  font=\sffamily,
  every node/.style={inner sep=0, outer sep=0},
  section-title/.style={font=\bfseries\Large, text=lighttext},
  section-kicker/.style={font=\small, text=mutetext},
  card-title/.style={font=\bfseries\normalsize, text=lighttext},
  card-sub/.style={font=\footnotesize, text=mutetext},
  body-text/.style={font=\small, text=lighttext, align=left},
  tiny-text/.style={font=\footnotesize, text=mutetext},
  arrow-flow/.style={-{Stealth[length=6pt]}, accent2, line width=1.2pt},
  arrow-data/.style={-{Stealth[length=5pt]}, accent1, line width=0.9pt, dashed},
]

% === BACKGROUND ===
\fill[darkbg] (0,0) rectangle (1200,1080);

% === HEADER ===
\node[section-title, anchor=west] at (60, 1025) {生成式艺术演示集};
\node[section-kicker, anchor=west] at (60, 998) {GENERA\kern0.7ptTIVE AR\kern0.7ptT DEMOS  /  12 个确定性 p5.js 透明 PNG 作品};

% Decorative line
\draw[primary, line width=2pt] (60,985) -- (360,985);

% Version badge
\node[fill=midbg, draw=bordercol, rounded corners=4pt, text=accent1, font=\footnotesize\bfseries, inner sep=4pt] at (1100, 1012) {v0.2.0};

% === SECTION 1: 项目概览 (left column) ===
\node[card-title, anchor=west] at (60, 950) {项目定位};
\draw[bordercol, line width=0.5pt] (60,939) -- (280,939);

\node[body-text, anchor=north west, text width=210pt] at (60, 924) {
一个以「完成态透明 PNG 构图」为目标的\\
生成式艺术参考库，而非持续动画。\\[6pt]
所有作品均使用固定随机种子（seed=42），\\
输出完全可复现。
};

% Key metrics row
\begin{scope}[shift={(60, 768)}]
  % Metric 1
  \fill[midbg] (0,0) rectangle (100,70);
  \draw[bordercol] (0,0) rectangle (100,70);
  \node[font=\bfseries\Large, text=primary] at (50, 42) {12};
  \node[font=\footnotesize, text=mutetext] at (50, 18) {艺术作品};

  % Metric 2
  \fill[midbg] (110,0) rectangle (210,70);
  \draw[bordercol] (110,0) rectangle (210,70);
  \node[font=\bfseries\Large, text=accent2] at (160, 42) {5};
  \node[font=\footnotesize, text=mutetext] at (160, 18) {主题色系};

  % Metric 3
  \fill[midbg] (220,0) rectangle (320,70);
  \draw[bordercol] (220,0) rectangle (320,70);
  \node[font=\bfseries\Large, text=accent1] at (270, 42) {1400};
  \node[font=\footnotesize, text=mutetext] at (270, 18) {输出宽度 / px};
\end{scope}

% === SECTION 2: 技术栈 ===
\node[card-title, anchor=west] at (60, 730) {核心技术栈};
\draw[bordercol, line width=0.5pt] (60,720) -- (280,720);

\begin{scope}[shift={(60, 697)}]
  % Tech items
  \foreach \name/\col/\y in {
    p5.js/primary/0,
    TypeScript/accent3/-35,
    Vite/accent2/-70,
    Playwright/accent4/-105,
    pngjs/accent1/-140
  } {
    \fill[\col!20] (0,\y-22) rectangle (8,\y-2);
    \node[body-text, anchor=west] at (16, \y-12) {\name};
  }

  % Right side descriptions
  \node[tiny-text, anchor=west, text width=130pt] at (130, -12) {Canvas 2D 渲染引擎};
  \node[tiny-text, anchor=west, text width=130pt] at (130, -47) {类型安全的源代码};
  \node[tiny-text, anchor=west, text width=130pt] at (130, -82) {开发服务器与构建};
  \node[tiny-text, anchor=west, text width=130pt] at (130, -117) {无头浏览器截图测试};
  \node[tiny-text, anchor=west, text width=130pt] at (130, -152) {PNG 像素级验证};
\end{scope}

% === SECTION 3: 测试验证 ===
\node[card-title, anchor=west] at (60, 498) {自动化验证};
\draw[bordercol, line width=0.5pt] (60,488) -- (280,488);

\node[body-text, anchor=north west, text width=220pt] at (60, 473) {
\textbf{\textcolor{accent2}{npm test}} 执行完整闭环：\\[4pt]
\begin{enumerate}
  \item[\textcolor{accent1}{1}] TypeScript 类型检查
  \item[\textcolor{accent1}{2}] Vite 构建产物
  \item[\textcolor{accent1}{3}] 启动本地服务器
  \item[\textcolor{accent1}{4}] Chrome 无头加载场景
  \item[\textcolor{accent1}{5}] 等待 \texttt{\_\_VIS\_READY\_\_}
  \item[\textcolor{accent1}{6}] 透明背景截图
  \item[\textcolor{accent1}{7}] 像素内容校验
  \item[\textcolor{accent1}{8}] 鼠标交互验证
\end{enumerate}
};

% Quality gate badge
\begin{scope}[shift={(60, 265)}]
  \fill[midbg] (0,0) rectangle (320,55);
  \draw[accent2!60, line width=1pt] (0,0) rectangle (320,55);
  \node[font=\bfseries\normalsize, text=accent2, anchor=west] at (15, 32) {质量门禁};
  \node[tiny-text, anchor=west, text width=280pt] at (15, 13) {
    透明度 $>8\%$ · 可见像素 $>3.5\%$ · 彩色像素 $>2500$
  };
\end{scope}

% === VERTICAL DIVIDER ===
\draw[bordercol, line width=0.8pt] (410,195) -- (410,930);

% === RIGHT SECTION: 12 件作品分类 ===
\node[card-title, anchor=west] at (440, 950) {作品家族 · 十二种生成算法};
\draw[bordercol, line width=0.5pt] (440,939) -- (1140,939);

% 4 families x 3 pieces grid
% Each piece: small icon box + title + tags

% Row 1: Particle & Field Systems
\begin{scope}[shift={(440, 908)}]
  \node[card-sub, font=\bfseries, text=accent2, anchor=west] at (0,0) {粒子与场系统 · PAR\kern0.7ptTICLE \& FIELD};
  \draw[accent2!40, line width=0.5pt] (0,-7) -- (700,-7);

  % Flow Field
  \begin{scope}[shift={(0,-50)}]
    \fill[midbg] (0,-44) rectangle (214, 24);
    \draw[bordercol] (0,-44) rectangle (214, 24);
    % Mini visual: dotted flow streamlines (non-crossing)
    \foreach \i in {0,1,2,3} {
      \foreach \j in {0,1,...,10} {
        \fill[primary!60] ({8+\j*8.2}, {-8-\i*8+2*sin(\j*36)}) circle (0.9pt);
      }
    }
    \node[font=\bfseries\small, text=lighttext, anchor=west] at (100, 8) {流场 Flow Field};
    \node[tiny-text, anchor=west] at (100, -8) {噪声粒子轨迹 · 220 粒种子};
    \node[tiny-text, anchor=north west, text width=104pt, align=left, text=primary!70] at (100, -15) {\#noise \#particles \#flow};
  \end{scope}

  % Reaction Diffusion
  \begin{scope}[shift={(232,-50)}]
    \fill[midbg] (0,-44) rectangle (214, 24);
    \draw[bordercol] (0,-44) rectangle (214, 24);
    % Mini visual: reaction dots
    \foreach \x in {8,16,...,88} {
      \foreach \y in {-8,-16,...,-40} {
        \pgfmathparse{int(random(0,100))}
        \ifnum\pgfmathresult>50
          \fill[accent2!70] (\x, \y) circle (1.5pt);
        \else
          \fill[accent1!50] (\x, \y) circle (1pt);
        \fi
      }
    }
    \node[font=\bfseries\small, text=lighttext, anchor=west] at (100, 8) {反应扩散 Reaction};
    \node[tiny-text, anchor=west] at (100, -8) {激活剂/抑制剂场 · 5 次迭代};
    \node[tiny-text, anchor=north west, text width=104pt, align=left, text=accent2!70] at (100, -15) {\#reaction-diffusion \#simulation};
  \end{scope}

  % Interference
  \begin{scope}[shift={(464,-50)}]
    \fill[midbg] (0,-44) rectangle (236, 24);
    \draw[bordercol] (0,-44) rectangle (236, 24);
    % Mini visual: one wave ring + off-centre source dots
    \draw[accent3!50, line width=0.5pt] (46, -12) circle (22pt);
    \fill[accent3!70] (46, -12) circle (2pt);
    \fill[accent4!60] (10, -34) circle (1.5pt);
    \fill[accent4!60] (84, -4) circle (1.5pt);
    \fill[accent4!60] (12, 10) circle (1.5pt);
    \fill[accent4!60] (82, -32) circle (1.5pt);
    \node[font=\bfseries\small, text=lighttext, anchor=west] at (100, 8) {干涉图谱 Interference};
    \node[tiny-text, anchor=west] at (100, -8) {5 个径向波源 · 相位合成};
    \node[tiny-text, anchor=north west, text width=126pt, align=left, text=accent3!70] at (100, -15) {\#waves \#interference \#field};
  \end{scope}
\end{scope}

% Row 2: Recursive & Dynamical Systems
\begin{scope}[shift={(440, 796)}]
  \node[card-sub, font=\bfseries, text=primary, anchor=west] at (0,0) {递归与动力系统 · RECURSIVE \& DYNAMICAL};
  \draw[primary!40, line width=0.5pt] (0,-7) -- (700,-7);

  % Botanical
  \begin{scope}[shift={(0,-50)}]
    \fill[midbg] (0,-44) rectangle (214, 24);
    \draw[bordercol] (0,-44) rectangle (214, 24);
    % Mini visual: L-system plant (filled silhouette, no strokes)
    \fill[accent2!45] (28,-38) rectangle (32,-18);
    \fill[accent2!55] (24,-13) circle (4.5pt);
    \fill[accent2!55] (36,-13) circle (4.5pt);
    \fill[accent2!55] (30,-22) circle (4.5pt);
    \fill[accent2!55] (30,-4) circle (4.5pt);
    \fill[primary!70] (24,-13) circle (1.3pt);
    \fill[primary!70] (36,-13) circle (1.3pt);
    \fill[primary!70] (30,-4) circle (1.3pt);
    \node[font=\bfseries\small, text=lighttext, anchor=west] at (100, 8) {植物递归 Botanical};
    \node[tiny-text, anchor=west] at (100, -8) {L 系统植物标本 · 6 层深度};
    \node[tiny-text, anchor=north west, text width=104pt, align=left, text=accent2!70] at (100, -15) {\#l-system \#recursion \#botanical};
  \end{scope}

  % Attractor
  \begin{scope}[shift={(232,-50)}]
    \fill[midbg] (0,-44) rectangle (214, 24);
    \draw[bordercol] (0,-44) rectangle (214, 24);
    % Mini visual: strange attractor scatter
    \foreach \i in {1,...,80} {
      \pgfmathsetmacro\x{50 + rnd*80 - 40}
      \pgfmathsetmacro\y{-20 + rnd*30 - 15 + sin(\i*37)*10}
      \fill[accent4!60] (\x pt, \y pt) circle (0.5pt);
    }
    \node[font=\bfseries\small, text=lighttext, anchor=west] at (100, 8) {奇异吸引子 Attractor};
    \node[tiny-text, anchor=west] at (100, -8) {Clifford Map · 24万次迭代};
    \node[tiny-text, anchor=north west, text width=104pt, align=left, text=accent4!70] at (100, -15) {\#attractor \#chaos \#trajectory};
  \end{scope}

  % Swarm
  \begin{scope}[shift={(464,-50)}]
    \fill[midbg] (0,-44) rectangle (236, 24);
    \draw[bordercol] (0,-44) rectangle (236, 24);
    % Mini visual: swarm boids ring
    \foreach \i in {0,10,...,350} {
      \pgfmathsetmacro\angle{\i}
      \pgfmathsetmacro\r{30 + mod(\i,60)*0.2}
      \pgfmathsetmacro\x{48 + cos(\angle)*\r}
      \pgfmathsetmacro\y{-18 + sin(\angle)*\r*0.5}
      \fill[accent1!70] (\x pt, \y pt) circle (0.8pt);
    }
    \node[font=\bfseries\small, text=lighttext, anchor=west] at (100, 8) {群体蜂拥 Swarm};
    \node[tiny-text, anchor=west] at (100, -8) {Boids  steering · 5 个族群};
    \node[tiny-text, anchor=north west, text width=126pt, align=left, text=accent1!70] at (100, -15) {\#boids \#agents \#motion};
  \end{scope}
\end{scope}

% Row 3: Spatial & Structural Systems
\begin{scope}[shift={(440, 684)}]
  \node[card-sub, font=\bfseries, text=accent1, anchor=west] at (0,0) {空间与结构系统 · SPA\kern0.7ptTIAL \& STRUCTURAL};
  \draw[accent1!40, line width=0.5pt] (0,-7) -- (700,-7);

  % Circle Packing
  \begin{scope}[shift={(0,-50)}]
    \fill[midbg] (0,-44) rectangle (214, 24);
    \draw[bordercol] (0,-44) rectangle (214, 24);
    % Mini visual: packed circles
    \fill[primary!50] (20, -14) circle (12pt);
    \fill[accent2!40] (48, -26) circle (8pt);
    \fill[accent1!50] (72, -12) circle (10pt);
    \fill[accent3!40] (38, -36) circle (5pt);
    \fill[accent4!50] (86, -30) circle (6pt);
    \fill[primary!30] (58, -6) circle (4pt);
    \node[font=\bfseries\small, text=lighttext, anchor=west] at (100, 8) {圆堆积 Circle Packing};
    \node[tiny-text, anchor=west] at (100, -8) {碰撞消解 · 150 个加权圆};
    \node[tiny-text, anchor=north west, text width=104pt, align=left, text=primary!70] at (100, -15) {\#packing \#collision \#composition};
  \end{scope}

  % Quadtree
  \begin{scope}[shift={(232,-50)}]
    \fill[midbg] (0,-44) rectangle (214, 24);
    \draw[bordercol] (0,-44) rectangle (214, 24);
    % Mini visual: quadtree subdivision (filled cells, gaps only)
    \fill[accent3!25] (8,-40) rectangle (46,-22);
    \fill[accent3!25] (52,-40) rectangle (90,-22);
    \fill[accent3!25] (8,-16) rectangle (46,2);
    \fill[accent3!45] (52,-16) rectangle (69,-7);
    \fill[accent3!45] (72,-16) rectangle (90,-7);
    \fill[accent3!45] (52,-4) rectangle (69,2);
    \fill[accent3!45] (72,-4) rectangle (90,2);
    \node[font=\bfseries\small, text=lighttext, anchor=west] at (100, 8) {四叉树 Quadtree};
    \node[tiny-text, anchor=west] at (100, -8) {自适应细分 · 6 层深度};
    \node[tiny-text, anchor=north west, text width=104pt, align=left, text=accent3!70] at (100, -15) {\#quadtree \#subdivision \#sampling};
  \end{scope}

  % Maze
  \begin{scope}[shift={(464,-50)}]
    \fill[midbg] (0,-44) rectangle (236, 24);
    \draw[bordercol] (0,-44) rectangle (236, 24);
    % Mini visual: maze wall segments (filled bars, gaps only)
    \fill[accent4!50] (8,-8) rectangle (30,-12);
    \fill[accent4!50] (38,-8) rectangle (60,-12);
    \fill[accent4!50] (66,-8) rectangle (88,-12);
    \fill[accent4!50] (18,-16) rectangle (22,-30);
    \fill[accent4!50] (48,-16) rectangle (52,-30);
    \fill[accent4!50] (74,-16) rectangle (78,-28);
    \fill[accent4!50] (8,-34) rectangle (36,-38);
    \fill[accent4!50] (44,-34) rectangle (70,-38);
    \fill[accent4!50] (76,-34) rectangle (88,-38);
    \node[font=\bfseries\small, text=lighttext, anchor=west] at (100, 8) {迷宫生成 Maze};
    \node[tiny-text, anchor=west] at (100, -8) {深度优先 · 39×21 网格};
    \node[tiny-text, anchor=north west, text width=126pt, align=left, text=accent4!70] at (100, -15) {\#maze \#algorithm \#paths};
  \end{scope}
\end{scope}

% Row 4: Pattern & Geometric Systems
\begin{scope}[shift={(440, 572)}]
  \node[card-sub, font=\bfseries, text=accent4, anchor=west] at (0,0) {图案与几何系统 · PA\kern0.7ptTTERN \& GEOMETRIC};
  \draw[accent4!40, line width=0.5pt] (0,-7) -- (700,-7);

  % Topography
  \begin{scope}[shift={(0,-50)}]
    \fill[midbg] (0,-44) rectangle (214, 24);
    \draw[bordercol] (0,-44) rectangle (214, 24);
    % Mini visual: dotted contour lines (non-crossing)
    \foreach \l in {0,1,...,4} {
      \foreach \j in {0,1,...,11} {
        \fill[accent2!55] ({7+\j*7.6}, {-36+\l*7+2*sin(\j*30)}) circle (0.8pt);
      }
    }
    \node[font=\bfseries\small, text=lighttext, anchor=west] at (100, 8) {地形图 T\kern0.7ptopography};
    \node[tiny-text, anchor=west] at (100, -8) {噪声等值线 · 42 层等高};
    \node[tiny-text, anchor=north west, text width=104pt, align=left, text=accent2!70] at (100, -15) {\#contour \#noise \#terrain};
  \end{scope}

  % Tapestry
  \begin{scope}[shift={(232,-50)}]
    \fill[midbg] (0,-44) rectangle (214, 24);
    \draw[bordercol] (0,-44) rectangle (214, 24);
    % Mini visual: woven tile checker (filled, no strokes)
    \foreach \x in {0,1,...,5} {
      \foreach \y in {0,1,2} {
        \pgfmathtruncatemacro\m{mod(\x+\y,2)}
        \ifnum\m=0
          \fill[primary!45] ({\x*14+8}, {\y*(-11)-6}) rectangle ({\x*14+18}, {\y*(-11)-14});
        \else
          \fill[accent1!40] ({\x*14+8}, {\y*(-11)-6}) rectangle ({\x*14+18}, {\y*(-11)-14});
        \fi
      }
    }
    \node[font=\bfseries\small, text=lighttext, anchor=west] at (100, 8) {挂毯纹样 T\kern0.7ptapestry};
    \node[tiny-text, anchor=west] at (100, -8) {24×12 纺织语法 · 模块组合};
    \node[tiny-text, anchor=north west, text width=104pt, align=left, text=primary!70] at (100, -15) {\#pattern \#tiles \#textile};
  \end{scope}

  % Tiling
  \begin{scope}[shift={(464,-50)}]
    \fill[midbg] (0,-44) rectangle (236, 24);
    \draw[bordercol] (0,-44) rectangle (236, 24);
    % Mini visual: fivefold ring of separated rhombi
    \foreach \i in {0,72,...,288} {
      \pgfmathsetmacro\cx{46+20*cos(\i)}
      \pgfmathsetmacro\cy{-14+14*sin(\i)}
      \fill[accent4!45] ({\cx+3.5}, \cy) -- (\cx, {\cy+2.5}) -- ({\cx-3.5}, \cy) -- (\cx, {\cy-2.5}) -- cycle;
    }
    \fill[accent4!70] (46,-14) circle (1.5pt);
    \node[font=\bfseries\small, text=lighttext, anchor=west] at (100, 8) {准晶平铺 Tiling};
    \node[tiny-text, anchor=west] at (100, -8) {五重对称菱形 · 11 环};
    \node[tiny-text, anchor=north west, text width=126pt, align=left, text=accent4!70] at (100, -15) {\#tiling \#geometry \#symmetry};
  \end{scope}
\end{scope}

% === ARCHITECTURE DIAGRAM: 渲染管线 ===
\node[card-title, anchor=west] at (440, 441) {渲染与测试管线};
\draw[bordercol, line width=0.5pt] (440,431) -- (1140,431);

% Pipeline diagram (fits 440..1140, rows end above the catalog section)
\begin{scope}[shift={(440, 409)}]
  % Box 1: Source Code
  \fill[midbg] (0,-70) rectangle (136, 10);
  \draw[accent3] (0,-70) rectangle (136, 10);
  \node[font=\bfseries\small, text=accent3] at (68, -10) {源代码};
  \node[tiny-text, text width=118pt] at (68, -35) {src/main.ts\\12 个渲染函数\\统一调色板};
  \node[tiny-text] at (68, -60) {p5.js + TypeScript};

  % Arrow 1->2
  \draw[arrow-flow] (136, -30) -- (188, -30);
  \node[tiny-text, text=accent2, anchor=south] at (162, -25) {Vite 构建};

  % Box 2: Vite Build
  \fill[midbg] (188,-70) rectangle (324, 10);
  \draw[accent2] (188,-70) rectangle (324, 10);
  \node[font=\bfseries\small, text=accent2] at (256, -10) {构建产物};
  \node[tiny-text, text width=118pt] at (256, -35) {dist/index.html\\静态资源\\可直接部署};
  \node[tiny-text] at (256, -60) {vite build};

  % Arrow 2->3
  \draw[arrow-flow] (324, -30) -- (376, -30);
  \node[tiny-text, text=accent2, anchor=south] at (350, -25) {本地服务};

  % Box 3: Dev Server
  \fill[midbg] (376,-70) rectangle (512, 10);
  \draw[primary] (376,-70) rectangle (512, 10);
  \node[font=\bfseries\small, text=primary] at (444, -10) {HTTP 服务};
  \node[tiny-text, text width=118pt] at (444, -35) {127.0.0.1:4174\\?scene= 参数切换\\离线渲染};
  \node[tiny-text] at (444, -60) {vite preview};

  % Arrow 3->4
  \draw[arrow-flow] (512, -30) -- (564, -30);
  \node[tiny-text, text=accent2, anchor=south] at (538, -25) {Playwright};

  % Box 4: Headless Chrome
  \fill[midbg] (564,-70) rectangle (700, 10);
  \draw[accent4] (564,-70) rectangle (700, 10);
  \node[font=\bfseries\small, text=accent4] at (632, -10) {无头浏览器};
  \node[tiny-text, text width=118pt] at (632, -35) {Chrome / SwiftShader\\等待 \_\_VIS\_READY\_\_\\鼠标交互验证};
  \node[tiny-text] at (632, -60) {playwright-core};

  % Arrow 4->5 (down)
  \draw[arrow-data] (632, -70) -- (632, -103);
  \node[tiny-text, text=accent1, anchor=west] at (638, -86) {截图输出};

  % Box 5: PNG Output (under box 4, serpentine flow)
  \fill[midbg] (252,-158) rectangle (700, -103);
  \draw[accent1] (252,-158) rectangle (700, -103);
  \node[font=\bfseries\small, text=accent1] at (476, -114) {透明 PNG 输出};
  \node[tiny-text, text width=420pt] at (476, -137) {out/*.png · 1400×900 · Alpha 通道保留 · omitBackground};

  % Arrow 5->6 (right-to-left)
  \draw[arrow-data] (252, -130) -- (188, -130);
  \node[tiny-text, text=accent1, anchor=south] at (220, -125) {像素校验};

  % Box 6: Validation
  \fill[midbg] (0,-158) rectangle (188, -103);
  \draw[accent2] (0,-158) rectangle (188, -103);
  \node[font=\bfseries\small, text=accent2] at (94, -114) {像素级验证};
  \node[tiny-text, text width=168pt] at (94, -137) {pngjs · 透明占比 · 可见像素 · 彩色像素阈值};
\end{scope}

% === CATALOG METADATA ===
\node[card-title, anchor=west] at (440, 219) {可检索目录};
\draw[bordercol, line width=0.5pt] (440,209) -- (1140,209);

\node[body-text, text width=680pt, anchor=north west] at (440, 195) {
  \texttt{catalog.json} 为每件作品定义标准化元数据，支持按\textbf{\textcolor{accent2}{用途}}、\textbf{\textcolor{primary}{家族}}、\textbf{\textcolor{accent1}{核心问题}}、\textbf{\textcolor{accent3}{复杂度}}和\textbf{\textcolor{accent4}{技术标签}}进行检索。
};

% Sample catalog entry
\fill[midbg] (620,116) rectangle (1140, 160);
\draw[bordercol] (620,116) rectangle (1140, 160);
\node[tiny-text, anchor=west, font=\ttfamily, text=accent2] at (635, 148) {
  \{ id: "flow-field", use: "particle systems", family: "particle systems",
};
\node[tiny-text, anchor=west, font=\ttfamily, text=mutetext] at (635, 136) {
  \quad question: "What trajectories emerge from a seeded curl-like field?",
};
\node[tiny-text, anchor=west, font=\ttfamily, text=mutetext] at (635, 124) {
  \quad tags: ["noise","particles","flow"] \}
};

% === FOOTER ===
\node[tiny-text, anchor=west] at (620, 24) {
  generative-art-demos · 确定性生成艺术参考库 · 固定种子可复现输出
};
\node[tiny-text, anchor=east] at (1140, 24) {
  \textcolor{primary}{/} \textcolor{accent1}{/} \textcolor{accent2}{/} \textcolor{accent3}{/} \textcolor{accent4}{/}  \quad 12 scenes · 5-color palette · seed: 42
};

\end{tikzpicture}
\end{document}
